\documentclass{article}

\usepackage[margin=1in,footskip=0.25in]{geometry}
\usepackage{setspace}
\doublespacing
\usepackage{booktabs}
\usepackage{hyperref}

\title{STROBE Checklist: Cross-Sectional Study}

\author{
Large language models for self-administered clinical vignette assessment:\\
a pilot and validation study in Vietnam\\[12pt]
Daniels BD, Zhang W(S), Nguyen H, Duong D
}

\date{}

\begin{document}

\maketitle

\noindent \textbf{Study Design:} Cross-sectional validation study

\noindent \textbf{Study Registration:} Not registered; study is a methods paper validating an assessment tool

\vspace{12pt}

\section*{STROBE Checklist for Cross-Sectional Studies}

This checklist indicates whether each of the 22 STROBE items is reported in the manuscript, and the corresponding page(s) or section(s) where information can be found.

\vspace{12pt}

\begin{table}[h]
\centering
\begin{tabular}{p{0.08\textwidth}p{0.65\textwidth}p{0.25\textwidth}}
\toprule
\textbf{Item} & \textbf{STROBE Checklist Item} & \textbf{Reported on} \\
\midrule

\textit{Title and abstract} & & \\
\midrule

1 & (a) Indicate the study's design with a commonly used term in the title or the abstract & Title, Abstract \\
  & (b) Provide in the abstract an informative and balanced summary of what was found & Abstract \\

\midrule
\textit{Introduction} & & \\
\midrule

2 & Give the rationale and objectives for this study, specifying any pre-specified hypotheses & Introduction, paragraphs 1--3 \\

\midrule
\textit{Methods} & & \\
\midrule

3 & Describe the key elements of the study setting (e.g., geographic location, number and type of institutions, in- or outpatient) & Materials and Methods, ``Study design and participants'' \\

4 & Describe the eligibility criteria, and sources and methods of participant selection and follow-up & Materials and Methods, ``Study design and participants'' \\

5 & (a) Give the number of participants at each stage of the study (e.g., screened, assessed for eligibility, approached, consented, analysed) & Materials and Methods, ``Study design and participants''; Results, Tables 1--2 \\
  & (b) Give reasons for non-participation at each stage & Materials and Methods, ``Study design and participants'' (minimal loss due to focused recruitment strategy) \\
  & (c) Consider use of a flow diagram & Not applicable; study is a recruitment design without substantial eligibility screening \\

6 & (a) Clearly define all outcomes, exposures, predictors, potential confounders, and effect modifiers. Give diagnostic criteria if applicable. & Materials and Methods, ``Outcomes'' and ``Data extraction and coding'' \\
  & (b) When applicable, describe sources of information and corresponding measurement dates. For repeated measures, describe the handling of repeat observations. & Materials and Methods, ``Data extraction and coding'' and ``Statistical analysis'' \\

7 & (a) Describe how quantitative variables were handled in the analyses. If applicable, describe which groupings were chosen and why & Materials and Methods, ``Statistical analysis''; Results \\
  & (b) If applicable, describe how missing data were addressed & Materials and Methods, ``Statistical analysis'' (no missing data issues reported) \\

8 & (a) Describe all statistical methods, including those used to control for confounding & Materials and Methods, ``Statistical analysis'' (Pearson correlation, OLS regression) \\
  & (b) Describe any methods used to examine subgroups and interactions & Materials and Methods, ``Statistical analysis'' (results calculated separately by scenario) \\
  & (c) Explain how matching of cases and controls was conducted, if applicable & Not applicable; not a case--control study \\
  & (d) Describe any sensitivity analyses & Discussion, paragraph 2 (future research directions; no sensitivity analyses performed in this study) \\

9 & Report numbers of outcome events or summary measures over time & Results, paragraphs 1--2; Tables 1--2; Figures 2--3 \\

10 & Describe characteristics of study participants (e.g., demographic, clinical, social) and information on exposures and potential confounders & Results, ``Focus group'' and ``Validation sample''; Tables 1--2 \\

11 & Give unadjusted and, if applicable, confounder-adjusted estimates and their precision (e.g., 95\% confidence interval). Make clear which confounders were adjusted for and why. & Results, paragraph 2 (Pearson correlations with 95\% confidence intervals); OLS regression coefficients and $r^{2}$ values \\

12 & Report category boundaries when continuous variables were categorized & Materials and Methods, ``Data extraction and coding'' (checklist items coded as binary; expert ratings on 1--5 scale) \\

\midrule
\textit{Results} & & \\
\midrule

13 & (a) Give characteristics of study participants (e.g., demographic, clinical, social) and information on exposures and potential confounders & Results, ``Focus group'' (9 providers) and ``Validation sample'' (22 providers, 132 interactions); Tables 1--2 \\
  & (b) Indicate the number of participants with missing data for each variable of interest & Results (no missing data reported) \\
  & (c) Summarise follow-up time and whether loss to follow-up occurred; reasons if applicable & Not applicable; cross-sectional study without follow-up \\

14 & (a) Report numbers of outcome events or summary measures & Results, ``Validation sample'' (Pearson correlations, regression coefficients) \\
  & (b) Report category boundaries when continuous variables were categorized & Materials and Methods, ``Data extraction and coding'' \\
  & (c) If applicable, estimate risk ratios, odds ratios, or differences in means and include 95\% confidence intervals and p-values & Results, ``Validation sample'' (Pearson $\rho$, OLS $r^{2}$, regression coefficients with confidence intervals) \\

15 & Present results of any investigations of interaction or subgroup analyses, including interaction or subgroup analyses specified a priori, as well as any subgroup analyses suggested by the data, with a note about whether they were prespecified & Materials and Methods, ``Statistical analysis'' (results calculated separately by scenario; not a hypothesis-testing study) \\

16 & Give results of all important analyses, including those not meeting the threshold for statistical significance. Present both unadjusted and adjusted results. & Results, ``Validation sample'' (all correlations and regression results reported); Figures 2--3 \\

\midrule
\textit{Discussion} & & \\
\midrule

17 & Summarize key results with reference to study objectives & Discussion, paragraph 1 \\

18 & Discuss limitations of the study, taking into account sources of potential bias or imprecision. Discuss both direction and magnitude of any potential bias. & Discussion, ``Limitations'' (sample size, generalizability, low LLM sensitivity, know-do gap, lack of non-verbal cues, LLM version dependency, connectivity requirements) \\

19 & Discuss the generalizability/external validity of the study findings & Discussion, paragraphs 1--2 and ``Limitations''; extended discussion of generalizability challenges and relevance to low- and middle-income country settings \\

20 & Give a general interpretation of the results in the context of other evidence & Discussion, paragraphs 1--3 (comparison to LLM-based virtual patient literature, clinical vignette methodology, cost implications) \\

\midrule
\textit{Other information} & & \\
\midrule

21 & Report funding sources and the role of funders. Report conflicts of interest for all authors. & Acknowledgments (Harvard Asia Center, Rose Service Learning, IQSS, Gilead Sciences) \\

22 & Provide information on how to access any supplementary information, code, and data made available & Materials and Methods, throughout; Supporting Information referenced (case prompts, platform code available open-source) \\

\bottomrule
\end{tabular}
\end{table}

\vspace{24pt}

\section*{Notes}

\noindent \textbf{Study Type:} This is a cross-sectional validation study assessing the performance of an LLM-based clinical vignette platform. The study includes two components: (1) a focus group usability testing phase ($n=9$), and (2) a validation phase ($n=22$ providers, 132 interactions).

\vspace{12pt}

\noindent \textbf{Key Deviations from Standard STROBE Format:}
\begin{itemize}
\item This is a methods/validation paper rather than a traditional epidemiological study. The primary goal is to validate a new assessment tool against expert clinical ratings, not to measure disease prevalence or investigate etiologic associations.
\item The study is not hypothesis-driven and does not employ statistical significance testing. Instead, effect sizes (correlations, $r^{2}$ values) are reported to quantify agreement between measures.
\item The study is reported following STROBE guidelines ``where applicable'' (as stated in the manuscript), with additional elements of the Standards for Reporting of Diagnostic Accuracy Studies (STARD) where relevant to the LLM-versus-human coding comparison.
\item No formal power calculation or sample size justification is provided, reflecting the exploratory nature of the validation exercise and the small sample size. However, the 132 interactions provide 433 unique data points.
\end{itemize}

\end{document}
